% Ad Atticum 10.13 --- headnote
% ad-atticum-10-13, 7 May 49 BC, Cumae

\textit{Cicero to Atticus, written from the Cuman villa
on the Nones of May 49 BC --- 7 May (the manuscript
dateline: \textit{Scr.\ in Cumano Non.\ Mai.\ a.\ 705
(49)}). Atticus's reply has arrived and lifted Tullia's
spirits, and Cicero's; he asks Atticus to keep writing
and not to let anything that points to hope slip past.
Then comes the sardonic centerpiece of the letter: do
not be frightened by Antony's lions (Antony was famous
for keeping a yoke of lions in his carriage at this
period); ``there is nothing more agreeable than the
man.'' Cicero proceeds to give Atticus a sample of a
``statesman's policy''
(\textit{[Greek: praxin politikou]}) at work: Antony
summoned the ten leading men and the boards of four
from the surrounding municipalities; they came at dawn
to his villa; Antony slept till the third hour; then,
when told the Neapolitans and Cumans were there (the
towns Caesar is angry with), he sent them away till the
next day --- he wanted a bath, and so on. Today he is
crossing to Aenaria (Ischia) to promise the exiles their
return. The mockery is bone-dry.}

\textit{Section 2 turns to housekeeping --- letters from
Axius, thanks for news of Tiro, the banker Vettienus,
the repayment to Vestorius, Servius Sulpicius's
expected arrival the next morning from Minturnae via
the Liternum estate of C.~Marcellus --- and to the
single thing that astonishes Cicero: Antony, who has
been particularly attentive in the past, has not even
sent him a messenger. Plainly the brief about Cicero is
harsher than expected; Antony does not want to refuse
him to his face. Section 3 returns to Spain --- the
hinge of everything --- and to Atticus's own
constraints (Silius, Ocella and the rest are held back;
even Atticus is held up by Curio). The closing word is
a Greek crux preserved daggered in the manuscripts as
\textit{[Greek: EKITAONON]} --- nonsense as it stands;
the editors mark it with $\dagger\dagger$ and leave it,
and so does this translation.}
